\section{Earlier Experiments: an Incomplete Archive, Restricted Menus, Controlled Refinements and Costs}\label{sec:exp}
This appendix retains the experiments of the previous version of the study. Their tasks, splits and summaries were specified before fitting their models.

\paragraph{A frozen comparison on preassigned PISA forms.}
The U.S. public-use student file of PISA 2012 contains 4,978 students in 162 schools \citep{pisaus2014}. The survey randomly assigned three questionnaire forms, each containing a common section and two of three rotating question sets, to keep estimated response time at 30 minutes \citep{pisa2014,kastberg2014}. We use the common report of grade repetition in grades 1--6 as the binary target and three binary response indicators from each of the three rotated question sets. Each student has six of the nine selected predictors, and none has all nine. Asked-but-unanswered items count as nonendorsements, while items absent from a student's form remain unobserved. The 226 records without a valid common target response are excluded. Target nonresponse is balanced across forms: 72, 75 and 79 nonrespondents on forms A, B and C (Pearson $\chi^2_2=0.33$, $p=0.85$); survey-weighted nonresponse rates are 4.4\%, 5.0\% and 5.2\%, with replicate-weight $z$-statistics of at most 1.3 for the pairwise differences.

Forty schools train the old ensemble, another forty train the new one, and the remaining 82 supply 2,441 audit records. Each ensemble averages three block-specific forests, fitted using only students who received that block. The contrast depends on all nine predictors. Every pair of blocks occurs in a form, so both the accuracy and Brier differences decompose into observable functions (Appendix~\ref{app:pisa}). Calibration uses incomplete training records only.

\begin{table}[t]
\centering
\caption{PISA secondary analysis on 2,441 held-out students. Brier improvement and 95\% intervals are multiplied by $10^3$; positive means the newer rule has smaller loss. Survey intervals use the 80 supplied Fay replicates. MI intervals also use Rubin combination under a working imputation model.}
\label{tab:pisa}
\small
\begin{tabular}{lrr}
\toprule
Method & Estimate & 95\% interval \\
\midrule
Form score & 1.13 & $[-2.28, 4.53]$ \\
Block/pair HT & 1.24 & $[-2.09, 4.57]$ \\
Logistic MI & 1.24 & $[-2.12, 4.60]$ \\
\bottomrule
\end{tabular}
\end{table}


Table~\ref{tab:pisa} compares optimized scores, a simple inverse-probability block/pair score, and twenty logistic chained-equation imputations with Rubin combination \citep{vanbuuren2011}. Survey-weighted intervals use the 80 supplied Fay replicates \citep{kastberg2014}. The audit domain is the set of students in the 82 audit schools; replicate estimates are computed for that domain with every replicate weight, the standard domain-estimation practice, and inference is conditional on the frozen models and the random school split. All three Brier intervals include zero, and the methods give similar estimates and widths. The randomized-accuracy gain is about 4.1 percentage points, a different target; both rules threshold to the negative class everywhere on this cube, so their deterministic accuracy difference is zero. The design interpretation assumes invariant reported answers across forms. Nine duplicated-item diagnostics find two unadjusted $p<0.05$ differences (ST29Q02, 4.0 points between forms A and B; ST42Q03, 4.8 points between forms B and C), neither surviving Bonferroni correction. As a sensitivity analysis we equalize each item's weighted endorsement rate across its two forms by random flips in the audit sample, lowering the higher form or raising the lower one, with 200 draws per scenario and the models and score tables held fixed. The mean shift of the optimized-score Brier estimate is at most 0.05 of its standard error, and the largest single shift is 0.35 standard errors. Equalizing both items at once shifts it by 0.015 standard errors on average. Complete-data truth and a branching counterfactual are unavailable, so this archive illustrates identification and uncertainty from real fixed forms, not an adaptivity ratio.

\paragraph{Independent forest updates with a restricted branching menu.}
On WDBC, Spambase and Ionosphere, twenty fixed splits each fit independent old and new random forests on disjoint training records \citep{wolberg1993,spambase1999,ionosphere1989}. Each forest has twenty depth-three trees. Nested sets of 3, 6, 9 and 12 training-binarized features give 240 tasks, all with cap $b=3$. The fixed menu contains every feature triple. The branching menu of the earlier analysis contained those forms and all forty component trees, with jointly optimized leaf scores.

\begin{table}[t]
\centering
\caption{Independent forests, cap 3: 60 binding tasks per dataset. Median fitted fixed/menu-branch variance ratio and oracle menu ratio (numerical lower endpoint). The branch menu includes every fixed triple and all 40 component trees. Neither column uses the unattainable full-forest target as its branching score.}
\label{tab:forest}
\small
\begin{tabular}{lrrr}
\toprule
Data & Fitted & Oracle menu & Fitted $>1.5$ \\
\midrule
WDBC & 1.085 & 1.090 & 0/60 \\
Spambase & 1.029 & 1.031 & 0/60 \\
Ionosphere & 1.117 & 1.100 & 0/60 \\
\bottomrule
\end{tabular}
\end{table}


With this menu the fitted variance ratios have medians 1.03--1.12 and none exceeds 1.5 (Table~\ref{tab:forest}, Figure~\ref{fig:forestbinding}). Appendix~\ref{app:forestexact} replaces the restricted menu by all depth-three trees. The exact ratios have median 1.51 over the 180 binding tasks, and the restricted menu captures a median of 17\% of the exact log-gain. Among 70 probability-output comparisons of the smaller independent-tree experiment, 60 already have active feature count above the cap, with population ratios of median 1.353 (SUSHI) and 1.295 (SAPA); for threshold outputs, 52 of 58 nonzero comparisons are nonbinding (Appendix~\ref{app:binding}).

\paragraph{Controlled gaps, transfer and operational costs.}
Depth-six CART/root refinements isolate a contrast that branching can compute in six queries. Exhaustive lower bounds over all fixed six-feature forms certify median population ratios of at least 1.883 on SUSHI and 1.474 on SAPA. The complete-calibration version is an efficiency diagnostic: its 500 independent target-population records already suffice for a direct audit at the chosen precision. Transferring calibration across SAPA age groups avoids that interpretation. Brier and clipped log-loss comparisons retain median fitted variance ratios 2.350 and 2.623, with conditional Wald coverage near 95\%. A separate restricted-pilot experiment charges all newly collected calibration records (Appendices~\ref{app:transfer} and~\ref{app:em}).

The imputation comparator in Tables~\ref{tab:fulldetail} and~\ref{tab:coverage} is a single deterministic regression imputation. Among the tree's path-set forms it selects the one whose calibration cell means best predict $Yd(X)$, asks that form of every audit record, and averages the calibration cell mean of $Yd(X)$ given the observed answers and label, with add-one shrinkage toward the overall calibration mean. It has no between-imputation variance component, so its intervals ignore calibration uncertainty, which explains its undercoverage. Proper multiple imputation is evaluated separately in Section~\ref{sec:experiments} and Appendix~\ref{app:screening}.

If the cap can be relaxed, compare complete acquisition using its actual label cost. With fixed/complete variance ratio $\rho>1$, fixed form size $b$ and complete size $m$, their continuous cost crossover is
\begin{equation}
 c_Y^*=\frac{m-b\rho}{\rho-1}.\label{eq:costcross}
\end{equation}
Median crossovers on the controlled tasks are 3.16 and 3.74 feature-cost units. At $c_Y=10$, complete records cost less on 45/45 SUSHI and 19/25 SAPA tasks (Figure~\ref{fig:costchoice}). When the cap is immutable, complete acquisition is inadmissible. For nonlinear updates, the bias certificate must be checked first: WDBC GBDT/MLP budget curves reject overly tight caps regardless of sample size (Appendix~\ref{app:budgetcurve}).

